\par\medskip\pdfbookmark[1]{Zu V.}{nexus-note-V}\centerline{Zu V.}\smallskip
Die Bemerkung zu Formel III zeigt, auf welchem Weg der Verfasser die Anzahl $k$ der in der Gesamtheit $C$ enthaltenen Formen bestimmt hat. Aus geometrischen Betrachtungen (vgl. I, Art. 5, und II, Art. 4) ergibt sich, dass der Grenzwert, dem sich der Quotient
\[
\frac{(1)+(2)+(3)+\ldots+(m)}m
\]
bei unbegrenzt wachsendem $m$ nähert, also die mittlere Anzahl der Darstellungen einer unbestimmten positiven ganzen Zahl,
\[
=k\frac\pi{\surd D}
\]
ist. Ein zweiter Ausdruck für denselben Grenzwert lässt sich auf verschiedene Arten aus der Beschaffenheit der zuvor bestimmten Anzahl $(m)=\epsilon f(m)$ der Darstellungen der Zahl $m$ ableiten. Nach den erhaltenen Bruchstücken (I, Art. 7, 8; III und IV) scheint der Verfasser dabei zunächst den folgenden Weg eingeschlagen zu haben.

Ist $\theta(m)$ eine beliebige Funktion der positiven ganzen Zahl $m$ und $p$ eine beliebige Primzahl, so lässt sich aus $\theta(m)$ stets eine neue Funktion $\theta'(m)$ ableiten, deren Wert unabhängig davon ist, ob und wie oft $p$ als Faktor in $m$ enthalten ist, und die für alle nicht durch $p$ teilbaren Zahlen $m$ mit $\theta(m)$ übereinstimmt. Man erhält eine solche Funktion, indem man $\theta'(m)=\theta\left(\frac m{p^\pi}\right)$ setzt, wobei $p^\pi$ die höchste in $m$ aufgehende Potenz von $p$ bezeichnet; man kann sagen, dass die Funktion $\theta'(m)$ aus $\theta(m)$ durch Elimination der Primzahl $p$ entsteht. Bildet man auf diese Weise aus $f(m)$ eine neue Funktion $f'(m)$ durch Elimination der Primzahl $2$, daraus die Funktion $f''(m)$ durch Elimination von $3$ und so weiter, so verläuft jede folgende Funktion regelmäßiger als die vorhergehenden. Eliminiert man die Primzahlen nacheinander in aufsteigender Reihenfolge, so hat eine solche Funktion
\SourceLeaf{305}
$\theta(m)$ für unendlich viele Werte von $m$ den Wert $f(1)=1$, insbesondere für alle Werte von $m$, die kleiner als die zuletzt eliminierte Primzahl sind. Durch unendliche Fortsetzung dieses Verfahrens nähert man sich immer mehr der Funktion $f^\infty(m)$, die für alle Werte von $m$ den Wert $1$ hat und deren mittlerer Wert folglich ebenfalls $=1$ ist. Gelingt es nun, den mittleren Wert einer Funktion $\theta(m)$ durch denjenigen der nächstfolgenden $\theta'(m)$ auszudrücken, so lässt sich auch der mittlere Wert der Funktion $f(m)$ durch eine unendliche Kette von Operationen finden.

Ist $p$ die Primzahl, durch deren Elimination $\theta'(m)$ aus $\theta(m)$ entsteht, so ist $\theta(m)=\theta'(m)f(p^\pi)$, wobei $p^\pi$ wieder die höchste in $m$ aufgehende Potenz von $p$ bezeichnet. Wenn $p$ kein Teiler von $2D$ ist, findet man daraus leicht
\[
\theta'(m)=\theta(mp)-\left(\frac{-D}p\right)\theta(m)
\]
Setzt man zur Abkürzung
\[
\begin{aligned}
\Theta(m)&=\theta(1)+\theta(2)+\ldots+\theta(m)\\
\Theta'(m)&=\theta'(1)+\theta'(2)+\ldots+\theta(m)
\end{aligned}
\]
so ergibt sich
\[
\Theta'(mp)=\Theta(mp)-\left(\frac{-D}p\right)\Theta(m)
\]
und daraus, wenn $\omega,\omega'$ die jeweiligen mittleren Werte der Funktionen $\theta(m),\theta'(m)$ bezeichnen,
\[
\omega=\frac{\omega'}{1-\left(\frac{-D}p\right)\frac1p}
\]

Wenn dagegen die Primzahl $p$ ein Teiler von $2D$ ist, besteht zwischen den Funktionen $\theta(m)$ und $\theta'(m)$ im Allgemeinen keine so einfache Beziehung mehr; auf ähnliche Weise ergibt sich jedoch leicht, dass in diesem Fall $\omega=\omega'$ ist. Die Beziehung zwischen $\omega$ und $\omega'$ lässt sich in beiden Fällen auch auf folgendem Weg ableiten. Setzt man
\[
\vartheta(m)=\Sigma\theta(\mu)
\]
wobei die Summe über alle Zahlen $\mu$ zu bilden ist, die nicht durch $p$ teilbar und zudem nicht größer als $m$ sind, und bezeichnet man mit $m',m'',m'''\ldots$ die jeweiligen größten in $\frac mp,\frac{m'}p,\frac{m''}p\ldots$ enthaltenen ganzen Zahlen, so ist
\begingroup\postdisplaypenalty=10000
\[
\begin{aligned}
\Theta(m)&=\vartheta(m)+\vartheta(m')f(p)+\vartheta(m'')f(pp)+\vartheta(m''')f(p^3)+\ldots\\
\Theta'(m)&=\vartheta(m)+\vartheta(m')+\vartheta(m'')+\vartheta(m''')+\ldots
\end{aligned}
\]
\endgroup
und daraus folgt
\[
\frac\omega{\omega'}=(1-\tfrac1p)\{1+\frac{f(p)}p+\frac{f(pp)}{pp}+\frac{f(p^3)}{p^8}+\ldots\}
\]
in Übereinstimmung mit dem oben gefundenen Ergebnis (vgl. die Note zu III und IV).

Der mittlere Wert der Funktion $f(m)$ ist daher gleich dem unendlichen Produkt
\[
\Pi\frac1{1-\left(\frac{-D}p\right)\frac1p}
\]
\SourceLeaf{306}
in dem $p$ alle Primzahlen durchlaufen muss, die keine Teiler von $2D$ sind; daraus folgt
\[
k=\frac{\epsilon\surd D}\pi\Pi\frac1{1-\left(\frac{-D}p\right)\frac1p}
\]

Hinsichtlich der Strenge dieser Herleitung bleibt jedoch ein Bedenken gegen die Methode, den mittleren Wert der Funktion $f(m)$ durch schrittweise Elimination aller Primzahlen zu bestimmen. Zwar ist ersichtlich, dass die durch Elimination der ersten $n$ Primzahlen erhaltene Funktion $f^{(n)}(m)$ im Wert mit der Funktion $f^\infty(m)=1$ übereinstimmt, solange $m$ kleiner als die zuletzt eliminierte Primzahl bleibt. Durch hinreichend große Wahl von $n$ lässt sich diese Übereinstimmung somit bis zu jeder vorgegebenen Größe der Zahl $m$ erreichen. Damit allein ist aber keineswegs bewiesen, dass sich bei unbegrenzt wachsendem $n$ der mittlere Wert von $f^{(n)}(m)$ dem mittleren Wert von $f^\infty(m)$, also dem Wert $1$, unbegrenzt nähert. Wie der Verfasser diese Lücke ausfüllen wollte, lässt sich aus den erhaltenen Papieren nicht mit Sicherheit erkennen. Doch führt die schon oben (in der Note zu I) mitgeteilte Formel
\[
1+4\{A-\frac A3+\frac A5-\frac A7+\frac A9-\frac A{11}\ldots\}
\]
für die Anzahl der Zahlenpaare, deren Quadratsumme den Wert $A$ nicht übersteigt, zu der Vermutung, dass der Verfasser unter Umgehung des unendlichen Produkts für den mittleren Wert der Funktion $f(m)$ unmittelbar die unendliche Reihe
\[
\Sigma\left(\frac{-D}n\right)\frac1n
\]
gefunden hat. Dabei muss $n$ der Größe nach alle positiven ganzen Zahlen durchlaufen, die zu $2D$ teilerfremd sind. Die einfachste Art, diesen Übergang anzudeuten, scheint die folgende zu sein.

\Needspace{6\baselineskip}
Ist $\mu$ der größte unter den Teilern einer Zahl $m$, die zu $2D$ teilerfremd sind, und setzt man $\theta(m)=f(\mu)$, so ist $\theta(m)$ die Funktion, die durch Elimination aller in $2D$ aufgehenden Primzahlen aus $f(m)$ entsteht und deren mittlerer Wert nach dem Obigen mit demjenigen von $f(m)$ übereinstimmt. Da nun $\theta(m)=\Sigma\left(\frac{-D}n\right)$ ist, wobei $n$ alle Teiler von $\mu$, also alle Teiler von $m$, die zu $2D$ teilerfremd sind, durchläuft, ergibt sich die zur obigen analoge Formel
\[
\Theta(m)=\theta(1)+\theta(2)+\ldots+\theta(m)=\Sigma\left(\frac{-D}n\right)\frac mn
\]
wobei in der rechten Summe der Buchstabe $n$ alle zu $2D$ teilerfremden Zahlen durchläuft und vom Quotienten $\frac mn$ stets nur die größte in ihm enthaltene ganze Zahl zu nehmen ist. Ordnet man die Glieder dieser Reihe nach wachsenden Zahlen $n$, so nimmt der Faktor $\frac mn$ fortwährend ab oder zumindest nie zu, und die Reihe bricht ab, sobald $n>m$ wird. Außerdem folgt aus dem Fundamentaltheorem der Theorie der quadratischen Reste und seiner Verallgemeinerung, dass die Summe von jeweils $\varphi(4D)$ aufeinanderfolgenden Werten des Faktors $\left(\frac{-D}n\right)$ verschwindet. Daraus ergibt sich, dass die Summe beliebig vieler aufeinanderfolgender Werte
\SourceLeaf{307}
dieses Faktors ihrem absoluten Wert nach niemals die endliche, nur vom Determinanten $D$ abhängige Größe $\Delta=\varphi(2D)$ übersteigt. Verbindet man diese beiden Bemerkungen, so findet man leicht, dass die Summe aller auf das Glied $\left(\frac{-D}n\right)\frac mn$ folgenden Glieder dem absoluten Wert nach kleiner als $\Delta\frac mn$ ist und folglich der Quotient $\Theta(m):m$ bei unendlich wachsendem $m$ die in der angegebenen Weise geordnete, konvergierende unendliche Reihe
\[
\Sigma\left(\frac{-D}n\right)\frac1n
\]
zum Grenzwert hat. Nachdem so der gemeinsame mittlere Wert der Funktionen $\theta(m)$ und $f(m)$ gefunden ist, erhält man unmittelbar
\[
k=\frac{\epsilon\surd D}\pi\Sigma\left(\frac{-D}n\right)\frac1n
\]

Ferner ist zu bemerken, dass die Artikel 6 und 8 der Abhandlung II auf eine in mancher Hinsicht einfachere und auch leicht durchführbare Behandlung des Problems hinweisen, bei der nur Darstellungen ungerader Zahlen oder sogar nur solcher Zahlen betrachtet werden, die zu $2D$ teilerfremd sind.
\BookRule{35mm}
